	\documentclass[10pt,conference]{IEEEtran}

\usepackage{cite}
\ifCLASSINFOpdf
   \usepackage[pdftex]{graphicx}
   \graphicspath{{figs/}}
   \DeclareGraphicsExtensions{.pdf,.jpeg,.png}
\else
   \usepackage[dvips]{graphicx}
   \graphicspath{{../figs/}}
   \DeclareGraphicsExtensions{.eps}
\fi

\usepackage[cmex10]{amsmath}
\interdisplaylinepenalty=2500
\usepackage{amsthm}
\newtheorem{definition}{Definition}
\usepackage{algorithmic}
\usepackage{array}
\usepackage{subcaption}
\usepackage{url}
\usepackage[T1]{fontenc}
\usepackage[utf8]{inputenc}
\usepackage[brazilian]{babel}
\usepackage{listings}

% corrija hifenação aqui
\hyphenation{op-tical net-works semi-conduc-tor}

\begin{document}
\title{Modelo de Relatório}

\newif\iffinal
\finalfalse
\finaltrue
\newcommand{\jemsid}{99999}

\iffinal
\author{\IEEEauthorblockN{Autor}
\IEEEauthorblockA{Bacharelado em Física \\
Instituto de Física\\
Universidade de Brasília \\
Email: aluno@aluno.unb.br}}
\fi


\maketitle

\begin{abstract}
Um resumo de 100-200 palavras. Deve conter, brevemente, motivação, problema estudado, solução proposta/implementada e resultados obtidos.
\end{abstract}
\IEEEpeerreviewmaketitle

\section{Introdução}

A introdução deve responder às seguintes perguntas de forma concisa e clara:

\begin{enumerate}
\item Qual é o problema estudado e o porquê de estudá-lo?
\item Qual a solução proposta e suas limitações?
\end{enumerate}

Na introdução é importante detalhar o problema proposto, quais as dificuldades em resolvê-lo a partir da solução proposta e mostrar de maneira concisa como será dada a análise do problema. Dica: apresentação, desenvolvimento, clímax, resolução. (Tamanho de texto: meia lauda).

\section{Metodologia}

Neste local do texto você deve descrever a ideia que foi usada para criar o código e não colocar o código em si. Equações matemáticas, fluxogramas e toda e qualquer forma de apresentar a metodologia utilizada deve estar contida nessa seção do relatório. (Tamanho de texto: meia lauda).

\section{Resultados e Discussão}

Descreva nesta seção os resultados encontrados de maneira detalhada e clara. Dica: ``do amplo para o específico'' ``mostrar e explicar'' (Tamanho de texto: uma lauda).

\section{Conclusão}

``So what? Statement'' (Tamanho de texto: meia lauda)

\end{document}


